\documentclass[12pt]{article}
\usepackage[T1]{fontenc}
\usepackage[utf8]{inputenc}
\usepackage{lmodern}
\usepackage{textcomp}
\usepackage{enumitem}
\usepackage{geometry}
\geometry{margin=1in}
\begin{document}
\begin{center}
Answer Key - Psychology - Graduate Level Assessment 4 \
\small (Answers are ordered exactly as in the question paper.)
\end{center}

\section*{Multiple Choice}
\begin{enumerate}[label=\arabic*.]
\item \textbf{Answer:} D
\\ \textit{Options:} A) tests that are used to measure intelligence; B) tests that are given only once a year; C) tests that have no right or wrong answers; D) tests for which a person's performance can be compared with a pilot group; E) any examination given by your state or country
\\ \textit{Note:} Standardized tests are designed to measure individual performance against a predetermined group, known as the pilot group, which establishes norms and benchmarks for interpreting test scores.
\item \textbf{Answer:} A
\\ \textit{Options:} A) The detection of spinal cord injuries; B) The detection of white matter abnormalities; C) The detection of soft tissue injuries; D) The detection of a small tumor; E) The detection of lung cancer
\\ \textit{Note:} MRI is not always preferred for detecting spinal cord injuries because CT scans are faster and more effective at visualizing acute traumatic injuries, making them more suitable in emergency settings.
\item \textbf{Answer:} A
\\ \textit{Options:} A) median; B) interquartile range; C) geometric mean; D) quartile deviation; E) variance
\\ \textit{Note:} In a left-skewed distribution, the median is positioned to the right of the mean, making it the highest value among central tendency measures, as it is less affected by extreme values in the tail.
\item \textbf{Answer:} A
\\ \textit{Options:} A) Kendal's tau; B) Biserial correlation; C) Cramer's V; D) Eta; E) Goodman and Kruskal's gamma
\\ \textit{Note:} Kendall's tau is the most appropriate correlation procedure for this situation because it is designed to assess the strength and direction of association between a dichotomous variable and an ordinal or continuous variable while being robust to non-normal distributions and outliers.
\item \textbf{Answer:} B
\\ \textit{Options:} A) extreme fatigue; B) high self-esteem; C) memory loss; D) intense fear and anxiety; E) uncontrollable grief and despair
\\ \textit{Note:} During the manic phase of bipolar disorder, individuals often exhibit a heightened sense of self-esteem and grandiosity, leading to an inflated self-image and unrealistic beliefs about their abilities and importance.
\end{enumerate}

\section*{True / False}
\begin{enumerate}[label=\arabic*.]
\setcounter{enumi}{5}
\item \textbf{Answer:} True
\\ \textit{Options:} A) True; B) False
\\ \textit{Note:} Episodic memory involves the ability to recall specific events and experiences, including those acquired within the recent past, thereby supporting the retention of information from the past few hours to days.
\item \textbf{Answer:} False
\\ \textit{Options:} A) True; B) False
\\ \textit{Note:} Anxiety disorders are influenced by both genetic predispositions and environmental factors, as research indicates that genetic heritability plays a significant role in their development.
\item \textbf{Answer:} False
\\ \textit{Options:} A) True; B) False
\\ \textit{Note:} The statement is false because proactive interference occurs when previously learned information (from list 1) disrupts the recall of new information (list 2), while list 3 does not contribute to this interference unless it is similar to list 2, thereby demonstrating that not all lists necessarily affect recall.
\item \textbf{Answer:} False
\\ \textit{Options:} A) True; B) False
\\ \textit{Note:} The statement is false because research indicates that individuals with varying levels of initial ability can demonstrate significant improvements in performance through extensive training, suggesting that effort and training quality can lead to growth regardless of starting skill level.
\item \textbf{Answer:} False
\\ \textit{Options:} A) True; B) False
\\ \textit{Note:} A learned response involves a behavior that is acquired through experience and conditioning rather than occurring automatically, which is characteristic of unlearned or reflexive responses.
\end{enumerate}

\section*{Long Form Response}
\begin{enumerate}[label=\arabic*.]
\setcounter{enumi}{10}
\item \textbf{Answer:} Client-centered therapy, developed by Carl Rogers, is fundamentally distinguished from classical therapeutic approaches by its foundational belief in human nature as innately capable and motivated towards self-fulfillment or actualization. This perspective posits that individuals possess an intrinsic drive to grow and develop, suggesting that the therapeutic relationship should be one of support rather than directive. In contrast, classical therapies often view individuals as needing guidance or correction, focusing on pathology and deficits rather than inherent potential. By emphasizing empathy, unconditional positive regard, and congruence, client-centered therapy creates an environment that fosters personal growth, allowing clients to explore and realize their true selves. This fundamental shift in perspective not only redefines the therapist-client dynamic but also empowers clients to take an active role in their healing and self- discovery.
\\ \textit{Note:} correct because it accurately identifies that client-centered therapy emphasizes the innate potential for self-fulfillment in individuals, contrasting with classical approaches that typically see clients as needing external guidance for change.
\item \textbf{Answer:} Case studies in qualitative research are in-depth examinations of specific instances or individuals, often aimed at exploring complex phenomena within their real-life context. A common misconception is that the findings from these case studies can be generalized to broader populations; however, this is misleading. Case studies prioritize the unique, contextualized insights gained from individual cases rather than seeking to produce findings that can universally apply. This focus on specificity contributes to theory development by providing rich, detailed data that can inform and refine existing theories or contribute to the emergence of new theoretical frameworks. Thus, while case studies may not yield generalizable results, they play a critical role in enhancing our understanding of psychological concepts and phenomena.
\\ \textit{Note:} correct because it emphasizes that case studies prioritize in-depth, context-specific insights over generalizability, highlighting their role in theory development rather than in producing universally applicable findings.
\end{enumerate}

\vfill
\noindent\textit{End of Answer Key}
\end{document}